\section{Generic NativeIPM vs. Parity-Specific IPM: Algorithm Gap and Refactoring Roadmap}
\label{sec:native-ipm-vs-parity-ipm}

This section isolates one question that now matters more than raw feature
completeness: why does the dedicated parity OPF interior-point solver converge
on large cases where the generic \texttt{NativeIPMAdapter} still fails on the
same nonlinear program callbacks?

The answer is no longer ``because the generic solver has no exact Hessian.''
That issue has already been addressed.  The remaining gap is that the dedicated
parity solver still has a better algorithmic envelope for OPF-class problems,
especially in globalization, centrality control, and reduced-system handling.

\subsection{What Is Actually Being Compared}

The large-case comparison in
\texttt{tests/benchmark\_ac\_opf\_cases.cpp} now evaluates three modes on the
same parity NLP model:
\begin{itemize}
  \item \texttt{parity}: the dedicated solver in \texttt{src/optimal_power_flow/parity\_ipm.cpp},
  \item \texttt{native}: the generic \texttt{NativeIPMAdapter} with the same
  callback model and exact Lagrangian Hessian,
  \item \texttt{ipopt}: IPOPT applied to the same callback model.
\end{itemize}

This matters because it removes most modeling ambiguity.  When \texttt{parity}
converges and \texttt{native} does not, the remaining explanation is mostly in
the solver algorithm, not in the OPF formulation itself.

\subsection{Algorithmic Differences}

Table~\ref{tab:native-vs-parity-ipm-gap} summarizes the most relevant current
differences.

\begin{longtable}{L{0.20\linewidth}L{0.34\linewidth}L{0.36\linewidth}}
\caption{Current algorithmic differences between generic NativeIPM and dedicated parity IPM}
\label{tab:native-vs-parity-ipm-gap}\\
\toprule
Aspect & Generic \texttt{NativeIPMAdapter} & Dedicated \texttt{parity\_ipm.cpp} \\
\midrule
\endfirsthead
\toprule
Aspect & Generic \texttt{NativeIPMAdapter} & Dedicated \texttt{parity\_ipm.cpp} \\
\midrule
\endhead
Globalization & Default path is filter-based with restoration and second-order correction; fallback path is merit-based backtracking.  In practice this stack is still fragile on large OPF NLPs. & Uses a more direct OPF-focused primal-dual step acceptance strategy with separate primal and dual fraction-to-boundary steps and no generic merit backtracking on each iteration. \\
Barrier-state initialization & Now improved with residual-aware slack and multiplier initialization in the merit path, but still embedded in a more generic driver. & Initializes slacks directly from inequality residuals to control \(\mu_i / z_i\) spread and maintain a well-conditioned condensed system from the start. \\
Residual metrics & Generic path now uses normalized residuals, but its acceptance logic still depends on a more generic merit/filter driver architecture. & Uses OPF-focused normalized convergence tests that were tuned around the parity formulation and its scaling behavior. \\
Complementarity handling & Uses a generalized predictor-corrector centering target and restoration logic. & Uses condensed complementarity handling with OPF-tested centering and direct update rules tailored to the assembled parity inequalities. \\
Extra correctors & Filter path contains second-order correction; merit path has no parity-style Gondzio block. & Includes explicit Gondzio-type extra corrector passes, which help recover centrality without relying on a merit backtracking loop. \\
Reduced Newton system & Reuses generic KKT utilities and multiple solve paths, including reduced-KKT logic for some equality patterns. & Builds and solves the condensed parity-specific Schur system directly in the exact structure expected by the OPF formulation. \\
Failure behavior on large parity OPF & Still terminates with globalization failure on IEEE118 and PEGASE2869. & Converges on IEEE118 and PEGASE2869 and passes PF cross-validation. \\
\bottomrule
\end{longtable}

\subsection{Why the Dedicated Parity Solver Still Wins}

The dedicated parity solver currently wins for three concrete reasons.

\paragraph{1. It solves the right condensed problem with less generic machinery}
The parity solver assembles the exact OPF inequality residuals, slack ratios,
and reduced Hessian contributions in one coherent path.  The generic solver can
now represent the same mathematics, but it still routes that information
through a broader abstraction that was designed to handle arbitrary NLPs rather
than one specific network structure.

\paragraph{2. Its globalization is still closer to the OPF central path}
The parity solver takes separate primal and dual fraction-to-boundary steps and
accepts them directly after corrector assembly.  The generic native solver,
especially in its merit-backed path, still has a stronger tendency to reject
steps via generic globalization rules that are mathematically sensible but too
conservative for the parity OPF trajectory.

\paragraph{3. It contains OPF-specific corrector behavior the generic solver still lacks}
The parity implementation includes Gondzio-style extra corrector passes and
best-iterate restoration in a path already tuned around parity OPF residual
geometry.  The generic solver has partially moved in this direction, but the
remaining gap is still visible on large cases.

\subsection{Current Performance Interpretation}

The present large-case evidence should be interpreted carefully.

\begin{itemize}
  \item The generic native NLP IPM is no longer catastrophically incomplete.  It
  now reaches much better neighborhoods on parity OPF than it did before the
  exact-Hessian and merit-path upgrades.
  \item However, the current performance is still not acceptable as a final OPF
  solver path.  On realistic parity cases it either fails outright or only
  approaches the correct trajectory before globalization collapses.
  \item Therefore the remaining issue is not ``can the solver form Newton
  systems?'' but rather ``can it keep and accept the right Newton trajectory on
  hard OPF landscapes?''
\end{itemize}

This is why the current Native IPM performance still looks problematic even
after recent improvements: the solver core is now mathematically much closer to
correct, but its driver logic is still not as OPF-effective as the parity
solver's dedicated update scheme.

\subsection{Implementation Backlog: PR1 / PR2 / PR3}

The next useful step is not another high-level roadmap but a concrete sequence
of pull requests.  The purpose of the following backlog is to bind the intended
algorithm changes directly to the current source functions and acceptance rules.

\paragraph{PR1: unify globalization around a direct primal/dual accepted step}

\textbf{Primary file:} \texttt{src/engine/ipm\_solver.cpp}

\textbf{Functions to change.}
\begin{itemize}
  \item \texttt{solve\_nlp\_filter\_impl}
  \item the merit fallback loop inside
  \texttt{NativeIPMAdapter::solve\_nlp\_detail}
  \item if needed, extract a shared helper for step acceptance so filter and
  merit paths do not diverge again
\end{itemize}

\textbf{What should change.}
\begin{itemize}
  \item Replace the current mixed acceptance rule in which \(x,s,\lambda\)
  are accepted by filter logic while \(\mu\) is still updated with an
  independent \texttt{alpha\_max\_dual}.  The accepted step should become one
  coherent primal-dual update decided from the same accepted trial.
  \item Introduce a direct fraction-to-boundary accepted step for the
  predictor-corrector direction before falling back to repeated generic
  backtracking.  In other words, the first-class path should be ``accept the
  Newton step if positivity and normalized residual improvement are good
  enough,'' not ``always reduce \(\alpha\) until the generic filter admits it.''
  \item Move the current acceptance decision away from pure
  \((\theta,\phi)\)-based rejection and toward a parity-style gate that checks
  positivity, normalized primal feasibility, normalized dual feasibility, and
  complementarity together.
\end{itemize}

\textbf{Acceptance criteria to replace or add.}
\begin{itemize}
  \item Replace ``accepted trial point for \(x,s,\lambda\) plus separate dual
  multiplier step'' with ``one accepted \((x,s,\lambda,\mu)\) trial.''
  \item Add an explicit rule that a full or near-full step is kept when it
  improves the normalized residual summary produced by
  \texttt{summarize\_residuals}, even if the generic filter would otherwise
  shrink it.
  \item Keep filter backtracking only as a secondary safeguard, not as the
  dominant accepted-step mechanism.
\end{itemize}

\textbf{Minimum validation.}
\begin{itemize}
  \item Re-run \texttt{benchmark\_ac\_opf\_cases} on IEEE118 and PEGASE2869 in
  \texttt{native} mode.
  \item Require that failure mode changes from early line-search collapse to a
  materially better trajectory, ideally either convergence or a much lower
  final normalized residual than the current baseline.
\end{itemize}

\paragraph{PR2: port parity-style extra correctors into the generic driver}

\textbf{Primary files:} \texttt{src/engine/ipm\_solver.cpp} and, for reference,
\texttt{src/optimal_power_flow/parity\_ipm.cpp}

\textbf{Functions to change.}
\begin{itemize}
  \item \texttt{solve\_nlp\_filter\_impl}
  \item the merit fallback loop inside
  \texttt{NativeIPMAdapter::solve\_nlp\_detail}
  \item possibly a new helper dedicated to extra-corrector assembly so the code
  does not duplicate one-off second-order logic in multiple branches
\end{itemize}

\textbf{What should change.}
\begin{itemize}
  \item Generalize the current one-shot second-order correction into a bounded
  parity-style extra-corrector loop.  The present filter path already has a
  single SOC retry; PR2 should turn that into a structured ``predictor,
  corrector, extra-corrector'' sequence.
  \item Use the extra correctors specifically to improve centrality and recover a
  larger accepted step, rather than treating them only as a feasibility patch
  after filter rejection.
  \item Ensure the extra corrector is evaluated against the same normalized
  residual summary used by the main step so that centrality improvement is
  measured rather than guessed.
\end{itemize}

\textbf{Acceptance criteria to replace or add.}
\begin{itemize}
  \item Replace ``at most one SOC attempt on first rejection'' with ``up to a
  small fixed number of extra correctors that are accepted only if residuals and
  complementarity improve.''
  \item Add an explicit centrality-improvement test based on the normalized
  complementarity and aggregate merit terms, not only on filter admissibility.
\end{itemize}

\textbf{Minimum validation.}
\begin{itemize}
  \item Re-run the same large-case parity benchmarks and compare accepted step
  lengths, iteration counts, and final residual summaries against PR1.
  \item Confirm that the generic solver is no longer losing obviously usable
  predictor-corrector steps that the parity solver would keep.
\end{itemize}

\paragraph{PR3: retune scaling, barrier-state initialization, and restoration triggers}

\textbf{Primary files:} \texttt{src/engine/ipm\_solver.cpp},
\texttt{include/hacdcpf/engine/ipm\_solver.hpp}, and any supporting scaling or
restoration helpers already used by the filter path

\textbf{Functions to change.}
\begin{itemize}
  \item \texttt{initialize\_barrier\_state}
  \item \texttt{summarize\_residuals}
  \item \texttt{solve\_nlp\_filter\_impl}
  \item restoration-trigger decisions inside the native driver stack
\end{itemize}

\textbf{What should change.}
\begin{itemize}
  \item Tighten the connection between slack initialization, multiplier
  initialization, and the normalized residual scales actually used during step
  acceptance.
  \item Revisit the current filter constants and restoration triggers so that the
  solver does not declare globalization failure too early on OPF-class problems
  that still have a promising best iterate.
  \item Make best-iterate restoration a deliberate part of the native path's
  acceptance policy rather than a final afterthought used only when the main
  loop is already exhausted.
\end{itemize}

\textbf{Acceptance criteria to replace or add.}
\begin{itemize}
  \item Replace purely absolute or mixed-scale progress decisions with
  consistently normalized criteria derived from \texttt{summarize\_residuals}.
  \item Add a restoration-entry rule based on repeated accepted-step collapse,
  not only on hard factorization or line-search failure.
  \item Keep the current best-iterate bookkeeping, but require restoration to be
  judged against parity-like feasibility and complementarity thresholds.
\end{itemize}

\textbf{Minimum validation.}
\begin{itemize}
  \item Re-run IEEE118 and PEGASE2869 in \texttt{native}, \texttt{parity}, and
  \texttt{ipopt} modes on the same callback model.
  \item Only declare the generic path ready for default use on parity OPF when
  it converges without IPOPT fallback, matches parity objective and feasibility
  closely, and no longer fails primarily through globalization breakdown.
\end{itemize}

\paragraph{Why this PR order matters}

This order is intentional.  PR1 attacks the most visible current defect, namely
that the accepted step is still too generic and internally inconsistent.  PR2
then adds the parity-style corrector machinery needed to keep that better step
on hard iterations.  PR3 only makes sense after PR1 and PR2, because scaling
and restoration thresholds cannot be tuned meaningfully while the accepted-step
logic itself is still unstable.

\subsection{Practical Bottom Line}

The current technical conclusion is blunt but useful.

\begin{itemize}
  \item The generic NativeIPM is now a genuine interior-point implementation and
  no longer a placeholder.
  \item Its current performance is still problematic on large OPF-class NLPs.
  \item The parity-specific solver remains the stronger implementation for this
  problem family.
  \item The right next move is not more Hessian plumbing, but better
  globalization and closer transfer of parity's accepted-step logic into the
  generic driver.
\end{itemize}

So if the review question is whether the current native NLP IPM still has major
performance problems, the answer is yes.  They are now concentrated in the
driver and globalization layer, not in the existence of the Newton system or
the availability of exact second derivatives.